\documentclass[11pt]{article}
\usepackage[margin=1in]{geometry}
\usepackage{amsmath,amssymb}
\title{Disproof of Conjecture 00000001464}
\author{TLMC AI agent submission}
\date{October 2026}
\begin{document}
\maketitle

\section{The conjecture}

\begin{quote}
\textit{Definition:} the Legendre transform $\nabla f^*$.
\textit{Conjecture:} beyond the period-2 property of the quadratic
Legendre iterate, there exists a convex function of period 3
(involution period three).
\end{quote}

\section{The refutation}

The Legendre transform $L$ on proper convex lower-semicontinuous
functions is an \emph{involution}: $L(L g) = g$ for all $g$
(Fenchel--Moreau).  Applying this to $g = Lf$ gives
\[
L\bigl(L(L f)\bigr) = L f \quad \text{for every } f .
\]
So if $f$ had period 3 --- $L(L(L f)) = f$ with $Lf \ne f$ --- the
identity would force $Lf = f$: period 3 collapses to period 1 (a
fixed point, e.g.\ the self-dual $f(x) = x^2/2$).  No convex function
of genuine period 3 exists; the period spectrum of $L$ is exactly
$\{1, 2\}$.

\section{Verification}

\texttt{reproduce.py} computes numerical Legendre transforms
(supporting lines on a fine grid) for convex examples: $x^2/2$ is a
fixed point; $|x|$, $\max(x^2/2, |x| - 1)$, and a Huber-like function
all satisfy $L(Lf) = f$ to machine zero --- period exactly 2, never 3.
The Lean package (core Lean~4, v4.33.1) treats $L$ abstractly with the
involution property as a hypothesis (Fenchel--Moreau, cited in prose)
and certifies \texttt{period3\_collapse},
\texttt{no\_genuine\_period3}, and the consistency of fixed points;
all four audited theorems report \texttt{does not depend on any
axioms}.

\section{Boundary}

The refutation applies to every instance of the abstract involution
axioms, in particular the concrete Legendre transform.  The numeric
round trips are on bounded grids with the supremum attained on-grid.

\end{document}
